\chapter{Avant-propos}

\section{À qui s'adresse ce cours}

À quelqu'un qui n'a jamais entraîné de modèle et qui voudrait le faire
\emph{sur sa propre machine}, sans louer de serveur et sans envoyer ses données
ailleurs.

Vous n'avez besoin ni de mathématiques avancées, ni d'expérience en
apprentissage automatique. Vous avez besoin d'un ordinateur avec une carte
graphique, d'un peu de patience, et de l'envie de comprendre ce que vous faites
plutôt que de recopier des commandes.

\section{Ce que vous saurez faire à la fin}

\begin{itemize}
  \item Expliquer, avec vos mots, ce que veut dire « entraîner un modèle ».
  \item Préparer votre machine, sur Windows, Linux ou macOS.
  \item Constituer un corpus — le texte dont le modèle va apprendre.
  \item Lancer un entraînement qui dure plusieurs jours et qui \textbf{survit}
        aux coupures de courant.
  \item Spécialiser un modèle existant sur votre domaine à vous.
  \item Dialoguer avec le modèle que vous venez d'entraîner.
  \item Savoir quel fichier contient quoi, et lequel donner à quelqu'un.
\end{itemize}

\section{Comment lire ce cours}

Les chapitres se suivent. Chacun s'appuie sur le précédent, et chacun se termine
par un exercice que vous pouvez faire en quelques minutes.

Trois encadrés reviennent tout au long du texte.

\begin{nkretenir}
Ce qu'il faut retenir : l'idée du passage, en deux ou trois phrases. Si vous ne
deviez lire qu'une chose par section, ce serait celle-là.
\end{nkretenir}

\begin{nkpiege}
Le piège : l'erreur qui coûte une soirée. Elles sont ici parce que quelqu'un
les a commises avant vous — souvent l'auteur de ce cours.
\end{nkpiege}

\begin{nkexo}
Exercice : à faire vraiment, pas à lire. On n'apprend pas à entraîner un modèle
en lisant, pas plus qu'on n'apprend à nager en regardant la piscine.
\end{nkexo}

\section{Une promesse, et une mise en garde}

\textbf{La promesse.} Toutes les commandes de ce cours ont été exécutées sur une
machine ordinaire — une carte graphique de 8~Go — avant d'être écrites ici. Les
chiffres que vous lirez (durées, tailles de fichiers, valeurs de perte) sont des
mesures, pas des estimations.

\textbf{La mise en garde.} Entraîner un modèle de sept milliards de paramètres
prend du \emph{temps}. Pas quelques minutes : plusieurs jours. Ce cours vous
apprendra à travailler avec cette contrainte plutôt que contre elle — c'est
d'ailleurs la raison d'être du chapitre~5, entièrement consacré à la reprise
après interruption.

\begin{nkretenir}
Vous n'avez pas besoin d'un ordinateur exceptionnel. Vous avez besoin de savoir
ce que vous faites, et d'accepter que la machine travaille pendant que vous
dormez.
\end{nkretenir}
